\begin{abstract}
Glutathione peroxidase~6 (GPX6) provides a natural system for asking how catalytic chemistry is reshaped by evolutionary sequence background. Primate GPX6 retains catalytic selenocysteine, whereas rodents have independently replaced it with cysteine, but the energetic consequences of this Sec/Cys exchange within each homologous scaffold remain unclear. Here, empirical valence bond simulations were used to quantify the proton-transfer step that generates the reactive selenolate or thiolate in human and mouse GPX6. Introducing selenocysteine into mouse GPX6 lowers the activation barrier by 2.84~kcal/mol, whereas replacing selenocysteine with cysteine in human GPX6 raises it by 1.43~kcal/mol, showing that the selenium advantage is amplified or damped by the surrounding protein scaffold. Greedy pathway free energy calculations with the empirical valence bond method across 20 proximal substitutions reveal multiple feasible routes between homologs and a recurring energetic pattern in which the effect of each substitution depends on the catalytic background in which it is introduced. Projection onto the GPX6 phylogeny indicates that rodent Cys-GPX6 is best understood as a compensatory catalytic regime rather than a simple loss of selenium chemistry.
\end{abstract}
